\cotefile{37}
\scanpage{31}
\margin{\figure{\begin{tikzpicture}\draw (0,0) -- (0,-6);\end{tikzpicture}}}
\uline{Théorème} 2.1.1. Soit $f:S'\handwritten{\longrightarrow}S$ un morphisme universellement localement \uline{$L$}-1-asphérique (SGA 4 XV \handwritten{1.11}), par exemple (SGA\handwritten{4}XV \handwritten{\unsure{2.2|2.1}}) un morphis\unsure{m} lisse, ou tel que $\deleted{\unsure{X|x|S}}\handwritten{S'}$ soit isomorphe à la limite\inserted{projective}\figure{\begin{tikzpicture}[x=1cm,y=1cm]\draw (0,0) .. controls (-.1,.35) and (.35,.65) .. (1,.65) -- (2.7,.65) .. controls (3,.65) and (3.2,.7) .. (3.2,1);\end{tikzpicture}} d'un systême projectif\inserted{filtrant}\figure{\begin{tikzpicture}[x=1cm,y=1cm]\draw (0,0) .. controls (-.1,.25) and (.2,.6) .. (.5,.65) -- (2.7,.95) -- (2.7,1.2);\end{tikzpicture}} de $S$-schémas lisses $S'_i$ à morphismes de transition affines\deleted{\unsure{,}} (SGA 4 XV \handwritten{1.13\inserted{(ii)}})\figure{\begin{tikzpicture}[x=1cm,y=1cm]\draw (0,0) .. controls (-.15,.5) and (-.15,.8) .. (.15,1) .. controls (.4,1.2) and (.4,1.35) .. (.3,1.5);\end{tikzpicture}}

Soi\inserted{en}t\figure{\begin{tikzpicture}[x=1cm,y=1cm]\draw (0,0) -- (.45,.8);\end{tikzpicture}} $g:X\handwritten{\longrightarrow}S$ un morphisme de schémas, d'où un carré car\unsure{\deleted{ét}té|\deleted{té}ét}sien
\margin{\figure{\begin{tikzpicture}\draw[line width=.3pt] (0,0) -- (.03,-.24);\draw[line width=.3pt] (.08,-.10) -- (.11,-.28);\end{tikzpicture}}}
\diagram{\begin{tikzcd}[ampersand replacement=\&]
X \handwritten{\arrow[d,"g"']} \& X' \handwritten{\arrow[l,"f'"']} \handwritten{\arrow[d,"g'"]} \\
S \& S'\handwritten{\arrow[l,"f"']}\rlap{,}
\end{tikzcd}}

$G$ un\inserted{ind-\uline{$L$}-}\figure{\begin{tikzpicture}[x=1cm,y=1cm]\draw (0,0) .. controls (-.05,.3) and (-.05,.65) .. (-.1,1) (0,.4) -- (1.8,.4) .. controls (2.1,.4) and (2.2,.55) .. (2.3,1.15);\end{tikzpicture}} champ sur $X$ \deleted{dont les \unsure{liux|liens|fiux} faisceaux d'automorphismes sont des ind-\uline{$L$}-faisceaux}\inserted{(\uline{V} \handwritten{5}\quad)}\unsure{,|\deleted{,}} $G'=f'^{\handwritten{*}}(G)$ le champ image inverse sur $X'$.

\begin{enumerate}[label=\alph*)]
\item Si \deleted{$S$ et $S'$ sont strictement locaux et} $f$ est un morphisme \deleted{strictement} \uline{local} de schémas strictement locaux, alors le foncteur
\[
G(X)\handwritten{\longrightarrow}G'(X')
\]
est une équivalence de catégories.

\item Si $g$ est cohérent(i.e. quasi-compact et quasi-séparé) alors le morphisme de champs ``de changement de base''
\[
f^{\handwritten{'}}(g_{\handwritten{*}}(G))\handwritten{\longrightarrow}g'_{\handwritten{*}}(f'^{\handwritten{*}}(G))\handwritten{\unsure{\deleted{=g_{s'}}|\deleted{=}g_{s'}|\deleted{=g_*}|\deleted{g_*}}}
\]
est une équivalence de champs sur $S'$.
\end{enumerate}

\margin{\figure{\begin{tikzpicture}\draw (0,0) -- (0,-2);\end{tikzpicture}}}
\uline{Corollaire} 2.1.2. Soient $k$ un corps séparablement clos, $k'$ une extension séparablement close de $k$, $X$ un $k$-schéma, $X'=X\handwritten{\otimes}_k k'$, alors pour tou\unsure{t} \deleted{champ $G$} ind-\uline{$L$}-champ $G$ sur $X$, désignant par $G'$ son image inverse sur $X'$, le foncteur
\[
G(X)\handwritten{\longrightarrow}G'(X')
\]
est une équivalence de catégories.

O\unsure{\deleted{q}n|\deleted{p}n|\deleted{n}q} sait en effet (SGA 4 XVI\quad\quad) que $\operatorname{Spec}(k')\handwritten{\longrightarrow}\operatorname{Spec}(k)$ est universellement localement \uline{$L$}-1-asphérique.

Démonstration. On prouve a) par référence à SGA 4 XVI (nouvelle édition !)\inserted{\handwritten{et à $\overline{\text{\uline{V}}}$ --}}\figure{\begin{tikzpicture}[x=1cm,y=1cm]\draw (0,0) .. controls (-.05,.3) and (.2,.5) .. (.6,.55) -- (2,.65) .. controls (2.45,.7) and (2.8,.85) .. (2.8,1);\end{tikzpicture}}, et b) en se ramenant à a) par des lemmes 2.1.7.2 et 2.1.7.4 renumérotés.
\margin{\handwritten{En tout, ça devrait faire 3 pages au Ch.p $\overline{\text{\uline{VII}}}$ au maximum, -- mais il y \unsure{aura|aurait} \unsure{renvois} au Ch.p $\overline{\text{\uline{\unsure{V|I}}}}$.}}
% archive sheet number: 30
